\providecommand{\toplevelprefix}{../..}  %
\documentclass[../../book-main_es.tex]{subfiles}

\begin{document}

\chapter*{Declaración de código abierto}

% \vspace{-2cm}
\begin{quote}
    «{\em A menudo comparo el código abierto con la ciencia. La ciencia
        ha adoptado toda esta noción del desarrollo abierto de ideas y de mejoras 
        en base a las ideas de otros. Esto hizo de la
        ciencia lo que es hoy en día y obtuvo los increíbles avances que hemos
        logrado. Y comparo eso con la brujería y la alquimia, donde la
        abertura era algo que no se practicaba}.»

    $~$\hfill -- Linus Torvalds
\end{quote}
\vspace{5mm}

Este libro está destinado a ser un proyecto de «código abierto» que se pondrá
a disposición del público en el sitio web
\begin{center}
    \url{https://ma-lab-berkeley.github.io/deep-representation-learning-book/}
\end{center}
para todos los interesadas en los principios científicos, matemáticos
y computacionales de la inteligencia. El libro se actualizará de manera
continua y periódica por todos los que quieran contribuir a
este proyecto, incluyendo sus versiones traducidas, los materiales educativos, 
y diversas herramientas de IA para apoyar
el estudio, la enseñanza y la investigación de este tema.

A través de este libro esperamos iniciar a un «programa de conocimiento 
abierto» que promueva que todo el conocimiento científico sea de código 
abierto, para su organización, pedagogía, difusión y avance.

\end{document}
